\documentclass[10pt,a4paper]{amsart}
\usepackage{amsmath,amssymb,amsthm,graphicx}
\usepackage{verbatim}
\newenvironment{reference}{\comment}{\endcomment}
%\newenvironment{reference}{}{}
\newenvironment{slogan}{\comment}{\endcomment}
\newenvironment{history}{\comment}{\endcomment}

% For commutative diagrams we use Xy-pic
\usepackage[all]{xy}

% We use 2cell for 2-commutative diagrams.
\xyoption{2cell}
% Theorem environments.
%
\theoremstyle{plain}
\newtheorem{theorem}[subsection]{Theorem}
\newtheorem{proposition}[subsection]{Proposition}
\newtheorem{lemma}[subsection]{Lemma}

\theoremstyle{definition}
\newtheorem{definition}[subsection]{Definition}
\newtheorem{example}[subsection]{Example}
\newtheorem{exercise}[subsection]{Exercise}
\newtheorem{situation}[subsection]{Situation}

\theoremstyle{remark}
\newtheorem{remark}[subsection]{Remark}
\newtheorem{remarks}[subsection]{Remarks}

\numberwithin{equation}{subsection}

% Macros
%
\def\lim{\mathop{\mathrm{lim}}\nolimits}
\def\colim{\mathop{\mathrm{colim}}\nolimits}
\def\Spec{\mathop{\mathrm{Spec}}}
\def\Hom{\mathop{\mathrm{Hom}}\nolimits}
\def\Ext{\mathop{\mathrm{Ext}}\nolimits}
\def\SheafHom{\mathop{\mathcal{H}\!\mathit{om}}\nolimits}
\def\SheafExt{\mathop{\mathcal{E}\!\mathit{xt}}\nolimits}
\def\Sch{\mathit{Sch}}
\def\Mor{\mathop{\mathrm{Mor}}\nolimits}
\def\Ob{\mathop{\mathrm{Ob}}\nolimits}
\def\Sh{\mathop{\mathit{Sh}}\nolimits}
\def\NL{\mathop{N\!L}\nolimits}
\def\CH{\mathop{\mathrm{CH}}\nolimits}
\def\proetale{{pro\text{-}\acute{e}tale}}
\def\etale{{\acute{e}tale}}
\def\QCoh{\mathit{QCoh}}
\def\Ker{\mathop{\mathrm{Ker}}}
\def\Im{\mathop{\mathrm{Im}}}
\def\Coker{\mathop{\mathrm{Coker}}}
\def\Coim{\mathop{\mathrm{Coim}}}

% Boxtimes
%
\DeclareMathSymbol{\boxtimes}{\mathbin}{AMSa}{"02}

%
% Macros for moduli stacks/spaces
%
\def\QCohstack{\mathcal{QC}\!\mathit{oh}}
\def\Cohstack{\mathcal{C}\!\mathit{oh}}
\def\Spacesstack{\mathcal{S}\!\mathit{paces}}
\def\Quotfunctor{\mathrm{Quot}}
\def\Hilbfunctor{\mathrm{Hilb}}
\def\Curvesstack{\mathcal{C}\!\mathit{urves}}
\def\Polarizedstack{\mathcal{P}\!\mathit{olarized}}
\def\Complexesstack{\mathcal{C}\!\mathit{omplexes}}
% \Pic is the operator that assigns to X its picard group, usage \Pic(X)
% \Picardstack_{X/B} denotes the Picard stack of X over B
% \Picardfunctor_{X/B} denotes the Picard functor of X over B
\def\Pic{\mathop{\mathrm{Pic}}\nolimits}
\def\Picardstack{\mathcal{P}\!\mathit{ic}}
\def\Picardfunctor{\mathrm{Pic}}
\def\Deformationcategory{\mathcal{D}\!\mathit{ef}}
\usepackage{cleveref}
\crefname{originallemma}{original lemma}{original lemmas}
\Crefname{originallemma}{Original lemma}{Original lemmas}
\makeatletter
\newlabel{lemma-limits-colimits-abelian-presheaves}{{03CL}{source}{}{sourcepointer-1}{}}
\newlabel{lemma-limits-colimits-abelian-presheaves@cref}{{[originallemma][0][]03CL}{[1][0][]source}{}{}{}}
\makeatother
\crefname{originalreference}{original reference}{original references}
\Crefname{originalreference}{Original reference}{Original references}
\makeatletter
\newlabel{sites-lemma-limit-sheaf}{{00W2}{source}{}{sourcepointer-2}{}}
\newlabel{sites-lemma-limit-sheaf@cref}{{[originalreference][0][]00W2}{[1][0][]source}{}{}{}}
\makeatother
\crefname{originalreference}{original reference}{original references}
\Crefname{originalreference}{Original reference}{Original references}
\makeatletter
\newlabel{categories-lemma-adjoint-exact}{{0038}{source}{}{sourcepointer-3}{}}
\newlabel{categories-lemma-adjoint-exact@cref}{{[originalreference][0][]0038}{[1][0][]source}{}{}{}}
\makeatother
\crefname{originalreference}{original reference}{original references}
\Crefname{originalreference}{Original reference}{Original references}
\makeatletter
\newlabel{homology-lemma-check-exactness}{{05AA}{source}{}{sourcepointer-4}{}}
\newlabel{homology-lemma-check-exactness@cref}{{[originalreference][0][]05AA}{[1][0][]source}{}{}{}}
\makeatother
\title{Limits and colimits of abelian sheaves on a site, and exactness of filtered colimits}
\author{The Stacks Project Authors}
\date{stacks-project:sites-modules.tex:local-snapshot-20261012; GNU FDL 1.2 or later}
\begin{document}
\maketitle
\noindent\textit{Editorial note.} All delivered mathematical text is the author's own, copied verbatim from the pinned Stacks Project chapter file sites-modules.tex; no statement or proof step was rewritten, reordered or removed.\par\medskip
\noindent\textit{Editorial note.} The counted span cites four external items (two lemmas and one categories lemma, plus the presheaf lemma), all recorded in assembly.external\_references with canonical tags as navigation pointers.\par\medskip
\noindent\textit{Allowed standard prerequisites (editorial).}\begin{itemize}
\item Limits and colimits of abelian presheaves are computed sectionwise. Stacks Project, Modules on Sites, Lemma lemma-limits-colimits-abelian-presheaves (tag 09B2). Supplies the presheaf-level descriptions that the sheaf-level statements are compared with.
\item The limit of a diagram of sheaves of sets is a sheaf; left adjoints preserve colimits. Stacks Project, Sites, Lemma sites-lemma-limit-sheaf (tag 00W1); Categories, Lemma categories-lemma-adjoint-exact (tag 0036). Justify that the presheaf limit is a sheaf and that sheafification computes the colimit, hence exactness for filtered colimits.
\end{itemize}\medskip
\ifdefined\hypertarget\hypertarget{sourcepointer-1}{}\fi
\noindent\textit{Source reference.} Original lemma 03CL: sources/stacks-chow/tags.txt, source lines 4388--4388. Cross-chapter standard prerequisite cited by the author inside the counted span; it is used as a pointer and its own proof is not part of this item.\par\medskip
\ifdefined\hypertarget\hypertarget{sourcepointer-2}{}\fi
\noindent\textit{Source reference.} Original reference 00W2: sources/stacks-chow/tags.txt, source lines 1168--1168. Cross-chapter standard prerequisite cited by the author inside the counted span; it is used as a pointer and its own proof is not part of this item.\par\medskip
\ifdefined\hypertarget\hypertarget{sourcepointer-3}{}\fi
\noindent\textit{Source reference.} Original reference 0038: sources/stacks-chow/tags.txt, source lines 125--125. Cross-chapter standard prerequisite cited by the author inside the counted span; it is used as a pointer and its own proof is not part of this item.\par\medskip
\ifdefined\hypertarget\hypertarget{sourcepointer-4}{}\fi
\noindent\textit{Source reference.} Original reference 05AA: sources/stacks-chow/tags.txt, source lines 6843--6843. Cross-chapter standard prerequisite cited by the author inside the counted span; it is used as a pointer and its own proof is not part of this item.\par\medskip
\begin{lemma}
\label{lemma-limits-colimits-abelian-sheaves}
Let $\mathcal{C}$ be a site.
\begin{enumerate}
\item All limits and colimits exist in $\textit{Ab}(\mathcal{C})$.
\item Limits are the same as the corresponding limits of abelian presheaves
over $\mathcal{C}$ (i.e., commute with taking sections over objects of
$\mathcal{C}$).
\item Finite direct sums are the same as the corresponding finite direct sums
in the category of abelian pre-sheaves over $\mathcal{C}$.
\item A colimit is the sheafification of the corresponding colimit in
the category of abelian presheaves.
\item Filtered colimits are exact.
\end{enumerate}
\end{lemma}

\begin{proof}
By Lemma \ref{lemma-limits-colimits-abelian-presheaves} limits and colimits
of abelian presheaves exist, and are described by taking limits and colimits
on the level of sections over objects.

\medskip\noindent
Let $\mathcal{I} \to \textit{Ab}(\mathcal{C})$, $i \mapsto \mathcal{F}_i$
be a diagram. Let $\lim_i \mathcal{F}_i$ be the limit of the diagram
as an abelian presheaf. By Sites, Lemma \ref{sites-lemma-limit-sheaf}
this is an abelian sheaf. Then it is quite easy to see that
$\lim_i \mathcal{F}_i$ is the limit of the diagram in
$\textit{Ab}(\mathcal{C})$. This proves limits exist and (2) holds.

\medskip\noindent
By Categories, Lemma \ref{categories-lemma-adjoint-exact}, and because
sheafification is left adjoint to the inclusion functor we see that
$\colim_i \mathcal{F}$ exists and is the sheafification of the colimit
in $\textit{PAb}(\mathcal{C})$. This proves colimits exist and (4) holds.

\medskip\noindent
Finite direct sums are the same thing as finite products in any abelian
category. Hence (3) follows from (2).

\medskip\noindent
Proof of (5). The statement means that given a system
$0 \to \mathcal{F}_i \to \mathcal{G}_i \to \mathcal{H}_i \to 0$
of exact sequences of abelian sheaves over a directed set $I$ the sequence
$0 \to \colim \mathcal{F}_i \to \colim \mathcal{G}_i \to
\colim \mathcal{H}_i \to 0$ is exact as well. A formal argument using
Homology, Lemma \ref{homology-lemma-check-exactness} and the
definition of colimits shows that the sequence
$\colim \mathcal{F}_i \to \colim \mathcal{G}_i \to \colim \mathcal{H}_i \to 0$
is exact. Note that $\colim \mathcal{F}_i \to \colim \mathcal{G}_i$
is the sheafification of the map of presheaf colimits which is
injective as each of the maps $\mathcal{F}_i \to \mathcal{G}_i$ is
injective. Since sheafification is exact we conclude.
\end{proof}

\end{document}
